\chapter*{Abstract}
\addcontentsline{toc}{chapter}{Abstract}

This thesis addresses EEG channel reconstruction under missing-channel conditions using a Graph Signal Processing (GSP) perspective, together with geometric baselines and temporal regularization methods. The central motivation is practical: missing channels in real EEG studies are rarely purely random and can severely distort both physiological interpretation and BCI downstream analysis.

Following the GSP view discussed by Ortega and collaborators, EEG is modeled as a graph signal. To ensure rigorous and fair evaluation, a comprehensive hyperparameter optimization pipeline was implemented (using the Optuna framework) across three distinct datasets (PhysioNet, BCI Competition IV 2a, and MNE Sample), simulating random and clustered (\textit{nearby}) missing channel scenarios.

The quantitative evaluation (measured by MAE, RMSE, DTW, and SNR) demonstrates a significant paradigm shift: temporal regularization methods, specifically TRSS and Tikhonov, substantially outperform geometric baselines, even when the latter are exhaustively optimized for best-case scenarios. Furthermore, TRSS excels at preserving physiological dynamics, successfully recovering Event-Related Potentials (ERP) and Power Spectral Density (PSD) under severe contiguous channel loss. These findings establish graph-temporal regularization as a robust, state-of-the-art framework for advanced EEG interpolation.
